\section{Saint Joseph de Maistre}

On \textbf{All Saint's Day}, we recall the guiding voices behind Gornahoor, most prominently: \textbf{Julius Evola}, \textbf{Rene Guenon}, \textbf{Vladimir Solovyov}, \textbf{Valentin Tomberg},
\textbf{August Comte}, \textbf{Donoso Cortes}. Then of course, there is our patron saint, \textbf{Joseph de Maistre}.

Joseph de Maistre was one of those well-born men of sound mind that Evola referred to in his self-defense. Although Evola appreciated Maistre more for his views opposing the French Revolution, he also reviewed Maistre's \textit{St Petersburg Dialogues}\footnote{\url{http://www.gornahoor.net/library/EvolaOnDeMaistre.pdf}} which we have translated.

Notice, in particular, the intelligent and cultured way that Evola deals with his theological differences with Maistre's Catholicism. When necessary, he addresses a particular issue of disagreement directly; otherwise, he seeks a broader framework, as in the discussion around original sin. He recognizes that what is essential is the common metaphysical view and the attitude of soul and spirit.

In another direction, we have \textbf{Rudolf Steiner}'s opinion that Maistre was ``a personality of the greatest imaginable genius, of compelling spirituality.'' Leaving aside Steiner's evolutionary worldview, he nevertheless recognizes a consciousness still existing in Europe. This is how Steiner summarizes Maistre:

\begin{quotex}
This view holds that since the beginning of the fifteenth century the course of human life on earth is going downhill. Since that time, only dissipation, godlessness, and vapidity have proliferated in European civilization; the mere intellect focusing on usefulness has gripped humanity. Truth, on the other hand, which is identical with the spirituality of the world, expresses something different since time immemorial. The problem is that modern man has forgotten this ancient, sacred truth. This primordial, sacred truth implies that man is a fallen creature. The human being has cause to appeal to his conscience and remorse in his soul so that he can lift himself up, so that his soul will not fall prey to materiality. But inasmuch as European humanity utilizes materiality since the middle of the fifteenth century, the European civilization is falling into ruin and with it the whole of mankind.
\end{quotex}


How astute, yet this could apply to Guenon, Evola and many others. Furthermore, anyone who does not recognize himself in that paragraph is \emph{ipso facto} on the other side. Steiner draws the battle line: his evolutionary, egalitarian perspective or that of Joseph de Maistre. This is a spiritual battle, which is not at all the same as an intellectual battle. Those who persist in viewing this struggle in terms of outward affiliations, such as paganism vs Christianity, or in terms of race or nation are missing the point, part of the underbrush that must be cleared away.

Joseph de Maistre, pray for us.


\flrightit{Posted on 2010-10-31 by Cologero}

\begin{center}* * *\end{center}

\begin{footnotesize}
\begin{sffamily}
\texttt{Lockheart on 2010-11-01 at 08:27 said:}

Amen.

\hfill

\texttt{Tosti on 2010-11-01 at 10:04 said:}

Nice and concise. Very good.

\hfill

\texttt{kadambari on 2010-11-20 at 11:16 said:}

Interesting perceptive article below. How the spirit of Russia was destroyed... \url{http://www.themoscowtimes.com/opinion/article/russia-could-have-been-china/422955.html}



\hfill

\end{sffamily}
\end{footnotesize}

